% =====================================================================
%  c05.tex —— §9 拓扑镜像对称（前半），p016–p021，(9.1)–(9.38)
% =====================================================================
\begin{exercise}
证明 $A[[\hbar]]$ 上 $\star$ 积的首项的对称部分 $B_1^+$ 可经规范变换化为零。
\end{exercise}

任意 Poisson 结构是否具有形变量子化（即是否存在以该 Poisson 结构为首项、定义在 $C^\infty(M)[[\hbar]]$ 上的 $\star$ 积）这一问题，已由 [Kon03] 肯定地解决。

\section{拓扑镜像对称}

Calabi--Yau 流形有如下若干互不等价的定义方式：
\begin{enumerate}
  \item 和乐群等于 $\mathrm{SU}(n)$ 的实 $2n$ 维 Riemann 流形；
  \item Ricci 曲率为 $0$ 的 Kähler 流形；
  \item 典范丛平凡的 Kähler 流形。
\end{enumerate}

\begin{exercise}
尽量详细地叙述上述三个定义之间的关系。
\end{exercise}

本节采用 3。此外，如无特别说明，假定 Calabi--Yau 流形是紧致的。由于典范丛平凡，处处不为零的全纯 $n$ 形式 $\Omega\in H^0(\Omega^n_Y)$ 在相差非零常数倍的意义下唯一存在，称之为全纯体积形式（holomorphic volume form）。

\begin{definition}
$n$ 维 Calabi--Yau 流形的组 $(Y,\check Y)$ 称为拓扑镜像对（topological mirror pair），是指对任意的 $0\le p,q\le n$ 均有
\begin{equation}
  h^{p,q}(Y)=h^{p,n-q}(\check Y)
  \tag{9.1}
\end{equation}
此时称 $\check Y$ 为 $Y$ 的拓扑镜像（topological mirror）。
\end{definition}

$Y$ 与 $\check Y$ 的 Hodge 数关于由 $q=n-q$ 所定义的直线呈镜像对称，这正是「镜像对称」一词的由来。
% [Tutor Note: 原文此处作「q = n − q で定義される直線」，按 Hodge 数的镜像对称，对称轴应为 q = n − p（即 p + q = n）；此处照原文保留。]

\begin{exercise}
设 $\check Y$ 是 $Y$ 的拓扑镜像，证明 $Y$ 与 $\check Y$ 的 Euler 数由
\begin{equation}
  \chi(Y)=(-1)^{n+1}\chi(\check Y)
  \tag{9.2}
\end{equation}
相联系。
\end{exercise}

\begin{remark}
如习题 1.18 所示，Hodge 数不是拓扑不变量。尽管如此，满足 (9.1) 的 Calabi--Yau 流形对仍被称为拓扑镜像对，一种解释是：当 $Y$ 是在满足 $H^1(\cO_X)\cong H^2(\cO_X)\cong 0$ 这一意义上的狭义 3 维 Calabi--Yau 流形时，非平凡的 Hodge 数只有 $h^{1,1}(X)$ 与 $h^{1,2}(X)$ 两个，而二者都由 Betti 数决定，即 $h^{1,1}(X)=b^2(X)$、$h^{1,2}(X)=b^3(X)/2-1$。
\end{remark}
% [Tutor Note: 原文此句前半用 $Y$，此处却写 $H^1(\cO_X)$、$h^{1,1}(X)$、$h^{1,2}(X)$，$X$ 与 $Y$ 混用，按原文保留。]

下述问题是拓扑镜像对称（topological mirror symmetry）中的中心课题：

\begin{problem}
对 Calabi--Yau 流形 $Y$，求其拓扑镜像存在的条件。
\end{problem}

已知具有给定 Chern 数 $c_1^2$ 与 $c_2$ 的一般型曲面的双有理等价类具有准射影的粗模空间 [Gie77]。对怎样的非负整数对 $(c_1^2,c_2)$ 模空间非空——追问这一点，便是一般型曲面的地理学（geography of surfaces of general type）[Per87] 的基本问题。这样的组只可能落在由宫冈--Yau 不等式 $c_1^2\le 3c_2$ [Miy77, Yau77] 与 Noether 不等式 $5c_1^2-c_2+36\ge 12h^{0,1}\ge 0$ 夹出的区域之内。容易看出，该区域中存在无穷多个使模空间非空的组 $(c_1^2,c_2)$。

拓扑镜像对称是关于 Calabi--Yau 流形的地理学的问题。至于 $n$ 维 Calabi--Yau 流形的 Hodge 数，连 $n=3$ 时模空间非空者是否有限都尚不清楚。

作为条件 (9.1) 的特殊情形，得到
\begin{equation}
  h^{1,1}(Y)=h^{1,n-1}(\check Y)
  \tag{9.3}
\end{equation}
而由与全纯体积形式的内积所确定的同构
\begin{equation}
  H^1(\Theta_Y)\xrightarrow{\ \sim\ }H^1(\Omega_Y^{n-1}),\qquad X\mapsto\iota_X\Omega
  \tag{9.4}
\end{equation}
可知 $h^{n-1,1}(\check Y)$ 与 $h^1(\Theta_{\check Y})$ 一致，后者再由小平--Spencer 理论与 Bogomolov--Tian--Todorov 定理而与复结构变形的自由度一致。另一方面，$h^{1,1}(Y)$ 是 $Y$ 的 Kähler 类的自由度的维数，因此 (9.3) 断言：$\check Y$ 的复结构变形的自由度与 $Y$ 的辛结构变形的自由度相等。这正是「镜像对称交换辛几何与复几何」这一原理的一处体现。

1 维的 Calabi--Yau 流形只有椭圆曲线，2 维的 Calabi--Yau 流形只有 Abel 曲面与 K3 曲面。椭圆曲线的 Hodge 菱形为
\begin{equation}
  \begin{array}{ccccc}
    & & 1 & & \\
    & 1 & & 1 & \\
    & & 1 & &
  \end{array}
  \tag{9.5}
\end{equation}
而 Abel 曲面与 K3 曲面的 Hodge 菱形分别为
\begin{equation}
  \begin{array}{ccccc}
    & & 1 & & \\
    & 2 & & 2 & \\
    1 & & 4 & & 1\\
    & 2 & & 2 & \\
    & & 1 & &
  \end{array}
  \tag{9.6}
\end{equation}
以及
\begin{equation}
  \begin{array}{ccccc}
    & & 1 & & \\
    & 0 & & 0 & \\
    1 & & 20 & & 1\\
    & 0 & & 0 & \\
    & & 1 & &
  \end{array}
  \tag{9.7}
\end{equation}
故它们全都是自对偶的。

\begin{remark}
复 2 维 Ricci 平坦 Kähler 流形除 Abel 曲面与 K3 曲面之外，还有 Enriques 曲面与双椭圆曲面（bi-elliptic surface）。前者是 K3 曲面关于非辛对合的商，后者是两个（未必同构的）椭圆曲线的直积关于有限群的商。
\end{remark}

\begin{exercise}
证明 Enriques 曲面的 Hodge 菱形由
\begin{equation}
  \begin{array}{ccccc}
    & & 1 & & \\
    & 0 & & 0 & \\
    0 & & 10 & & 0\\
    & 0 & & 0 & \\
    & & 1 & &
  \end{array}
  \tag{9.8}
\end{equation}
给出。
\end{exercise}

\begin{exercise}
证明双椭圆曲面的 Hodge 菱形由
\begin{equation}
  \begin{array}{ccccc}
    & & 1 & & \\
    & 1 & & 1 & \\
    0 & & 2 & & 0\\
    & 1 & & 1 & \\
    & & 1 & &
  \end{array}
  \tag{9.9}
\end{equation}
给出。
\end{exercise}

K3 曲面与 Enriques 曲面的和乐群是 $\mathrm{SU}(2)$，而 Abel 曲面与双椭圆曲面的和乐群是平凡群。

3 维 Calabi--Yau 流形中大概最有名的例子是 $\PP^4$ 的 5 次超曲面。把它记作 $Y$，其定义方程由 5 个变数的 5 次式给出。典范丛平凡一事，可由对正合列
\begin{equation}
  0\to T_Y\to T_{\PP^4}|_Y\to N_{Y/\PP^4}\to 0
  \tag{9.10}
\end{equation}
取最高次外积所得的伴随公式（adjunction formula）
\begin{equation}
  \omega_Y^\vee\otimes N_{Y/\PP^4}\cong\omega_{\PP^4}^\vee|_Y
  \tag{9.11}
\end{equation}
直接看出。由 Lefschetz 定理，$Y$ 的低于中间维数的 Hodge 数与 $\PP^4$ 的 Hodge 数相同，从而 $h^{1,0}(Y)=h^{2,0}(Y)=0$ 且 $h^{1,1}(Y)=1$。由典范丛平凡可知 $h^{3,0}(Y)=1$。一般地，据 Griffiths [Gri69]，由次数为 $h$ 的 $n+2$ 元齐次多项式 $f$ 定义的 $\PP^{n+1}$ 的超曲面的中间维数 Hodge 数由
\begin{equation}
  h^{p,n-p}=\dim\operatorname{Jac}(f)_{h(n+1-p)-(n+2)}
  \tag{9.12}
\end{equation}
给出。其中
\begin{equation}
  \operatorname{Jac}(f):=\C[x_1,\ldots,x_{n+2}]\Big/\left(\frac{\partial f}{\partial x_1},\cdots,\frac{\partial f}{\partial x_{n+2}}\right)
  \tag{9.13}
\end{equation}
是 $f$ 的 Jacobi 环。

\begin{exercise}
证明 $\PP^4$ 的 5 次超曲面的 Hodge 菱形由
\begin{equation}
  \begin{array}{ccccccc}
    & & & 1 & & & \\
    & & 0 & & 0 & & \\
    & 0 & & 1 & & 0 & \\
    1 & & 101 & & 101 & & 1\\
    & 0 & & 1 & & 0 & \\
    & & 0 & & 0 & & \\
    & & & 1 & & &
  \end{array}
  \tag{9.14}
\end{equation}
给出。
\end{exercise}

\begin{exercise}
证明 $\PP^4$ 的 5 次超曲面即使稍微变形复结构，也仍保持为 $\PP^4$ 的射影超曲面。
\end{exercise}
% [Tutor Note: 原文「…のままであることとを示せ」多出一个「と」，疑为误植。]

\begin{exercise}
证明 $\PP^4$ 的 5 次超曲面的任意全纯自同构都提升为 $\PP^4$ 的自同构。
\end{exercise}

\begin{exercise}
利用习题 9.6 与习题 9.7，证明 $\PP^4$ 的 5 次超曲面的复结构模空间的维数为 101。
\end{exercise}

$\PP^4$ 的 5 次超曲面的拓扑镜像 $\check Y$ 由 [GP90] 如下构造。首先，作为 $\PP^4$ 的 5 次超曲面的特别例子，取 Fermat 流形 $Y:=\left\{[x_1:\cdots:x_5]\in\PP^4\;\middle|\;x_1^5+x_2^5+x_3^5+x_4^5+x_5^5=0\right\}$。这个 $Y$ 上自然地作用着 $G:=\left\{\operatorname{diag}(\alpha_1,\ldots,\alpha_5)\in \mathrm{SL}_5(\C)\;\middle|\;\alpha_1^5=\cdots=\alpha_5^5=1\right\}\cong(\Z/5\Z)^4$，商空间 $Y/G$ 是带商奇点的代数簇。该代数簇具有 crepant 解消，故把它记作 $\check Y$。

\begin{exercise}
计算 $\check Y$ 的 Hodge 数，并验证 $\check Y$ 是 $Y$ 的拓扑镜像。
\end{exercise}

crepant 解消并不唯一，但它们由 flop 相连，且 Hodge 数与 crepant 解消的选取无关 [Kol89]。关于具体的奇点解消的描述，例如参见 [Mor93, Appendix B] 及其参考文献。

3 维 Calabi--Yau 流形还有如下例子：椭圆曲线的直积
\begin{equation}
  \begin{aligned}[b]
    E^3:=\big\{&([x_1:x_2:x_3],[x_4:x_5:x_6],[x_7:x_8:x_9])\in(\PP^2)^3\mid{}\\
    &x_1^3+x_2^3+x_3^3=x_4^3+x_5^3+x_6^3=x_7^3+x_8^3+x_9^3=0\big\}
  \end{aligned}
  \tag{9.15}
\end{equation}
上作用着 $G:=\left\{(\alpha,\beta)\in(\C^\times)^2\;\middle|\;\alpha^3=\beta^3=1\right\}$，作用为
\begin{equation}
  \begin{aligned}[b]
    &([x_1:x_2:x_3],[x_4:x_5:x_6],[x_7:x_8:x_9])\\
    &\qquad\mapsto([\alpha\beta x_1:x_2:x_3],[\alpha^2x_4:x_5:x_6],[\beta^2x_7:x_8:x_9])
  \end{aligned}
  \tag{9.16}
\end{equation}
商流形 $E^3/G$ 具有 crepant 解消
\begin{equation}
  Y\to E^3/G
  \tag{9.17}
\end{equation}

\begin{exercise}
证明 $Y$ 的 Hodge 菱形由
\begin{equation}
  \begin{array}{ccccccc}
    & & & 1 & & & \\
    & & 0 & & 0 & & \\
    & 0 & & 84 & & 0 & \\
    1 & & 0 & & 0 & & 1\\
    & 0 & & 84 & & 0 & \\
    & & 0 & & 0 & & \\
    & & & 1 & & &
  \end{array}
  \tag{9.18}
\end{equation}
给出。
\end{exercise}

由于 Kähler 流形的 $h^{1,1}$ 必为正，$Y$ 在定义 9.1 的意义下没有拓扑镜像。一般地，满足 $h^{1,n-1}(Y)=0$ 的 Calabi--Yau 流形因没有变形而被称为刚性 Calabi--Yau 流形（rigid Calabi--Yau manifold）。关于刚性 Calabi--Yau 流形的镜像，后面还会稍作叙述。

作为 $\PP^4$ 的 5 次超曲面的推广，有加权射影空间的超曲面。正整数的组 $(a_1,\ldots,a_{n+2})$ 称为权（weight）。给定权后，乘法群 $\mathbb{G}_m$ 通过同态
\begin{equation}
  \mathbb{G}_m\to \mathrm{GL}_{n+2},\qquad \alpha\mapsto\operatorname{diag}(\alpha^{a_1},\ldots,\alpha^{a_{n+2}})
  \tag{9.19}
\end{equation}
作用在 $\mathbb{A}^{n+2}$ 上；该作用关于 $\mathbb{A}^{n+2}\setminus 0$ 的商称为加权射影空间（weighted projective space），记作 $\PP(a_1,\ldots,a_{n+2})$。加权射影空间是至多只有商奇点的代数簇。对任意正整数 $d$，把所有权统一乘以 $d$ 所得的加权射影空间同构：
\begin{equation}
  \PP(da_1,\ldots,da_{n+2})\cong\PP(a_1,\ldots,a_{n+2}).
  \tag{9.20}
\end{equation}

又注意到 $\mathbb{A}^1/(\Z/d\Z)\cong\mathbb{A}^1$，当 $\gcd\left((a_i)_{i=1}^{n+2}\right)=1$ 时还可看出：把除第 $i$ 个以外的权统一除以 $\gcd\left((a_j)_{j\ne i}\right)$，所得的加权射影空间仍同构于原来的加权射影空间。因此，可以假定对任意 $i=1,\ldots,n+2$ 都有
\begin{equation}
  \gcd\left(\{a_1,\ldots,a_{n+2}\}\setminus\{a_i\}\right)=1
  \tag{9.21}
\end{equation}
满足该条件的权称为良构（well-formed）的。$n+2$ 元多项式 $f(x_1,\ldots,x_{n+2})$ 若对任意 $\alpha\in\mathbb{G}_m$ 都满足 $f(\alpha^{a_1}x_1,\ldots,\alpha^{a_{n+2}}x_{n+2})=\alpha^hf(x_1,\ldots,x_{n+2})$，则称为 $h$ 次加权齐次多项式（weighted homogeneous polynomial）。

令
\begin{equation}
  \tilde Y:=\left\{(x_1,\ldots,x_{n+2})\in\mathbb{A}^{n+2}\;\middle|\;f(x_1,\ldots,x_{n+2})=0\right\}
  \tag{9.22}
\end{equation}
则
\begin{equation}
  Y:=\left(\tilde Y\setminus 0\right)/\mathbb{G}_m
  \tag{9.23}
\end{equation}
确定 $\PP(a_1,\ldots,a_{n+2})$ 的超曲面。当 $\tilde Y$ 除原点外没有奇点时，称 $Y$ 是拟光滑（quasi-smooth）的。拟光滑超曲面只具有来自作为容器的加权射影空间的商奇点。当 $Y$ 是反典范除子（anti-canonical divisor）时，即当 $Y$ 的定义方程的次数 $h$ 满足
\begin{equation}
  a_1+\cdots+a_{n+2}=h
  \tag{9.24}
\end{equation}
时，由伴随公式可知 $Y$ 具有平凡的典范丛。但若取任意的权 $(a_1,\ldots,a_{n+2})$，则无论怎样取次数为 $a_1+\cdots+a_{n+2}$ 的齐次多项式，$Y$ 都不会拟光滑。例如取 $n=1$、$(a_1,a_2,a_3)=(1,1,3)$，一般的 5 次多项式 $f(x_1,x_2,x_3)$ 可用 $x_1$ 与 $x_2$ 的 5 次式 $g(x_1,x_2)$ 及 2 次式 $h(x_1,x_2)$ 写成
\begin{equation}
  f(x_1,x_2,x_3)=g(x_1,x_2)+h(x_1,x_2)x_3
  \tag{9.25}
\end{equation}
因而必定在 $(x_1,x_2,x_3)=(0,0,1)$ 处具有奇点。

使一般的反典范除子拟光滑的权的集合，在 $n=1$ 时只有 $(1,1,1)$、$(1,1,2)$、$(1,2,3)$ 三个 [Sai74]。相应的加权射影平面的超曲面是椭圆曲线，对它们而言拓扑镜像对称没有深刻内容。

$n=2$ 的情形已由 [Yon90] 分类，已知共有 95 个。相应的加权射影空间超曲面的极小模型给出 K3 曲面，但与椭圆曲线一样，K3 曲面是自对偶的，因此拓扑镜像对称照原样并无深刻内容。然而，若注意到这些 K3 曲面作为族存在，便可把对它们的拓扑镜像对称重新表述为如下对偶性：非常一般（very general）\footnote{给定 $\C$ 上的代数簇 $V$ 时，若某性质在从 $V$ 中去掉至多可数个非平凡的 Zariski 闭集所得的集合上成立，则称该性质在 $V$ 的非常一般的点（very general point）处成立。}的元素的 Picard 格与超越格互换 [Dol96]。这与 Arnold [Arn75] 发现、并由 Pinkham [Pin77] 及 Dolgachev--Nikulin [Dol83, Nik79] 用 K3 曲面的几何加以解释的奇妙对偶性（strange duality）密切相关\footnote{关于 K3 曲面的镜像对称，另见 [Ued] 及其参考文献。}。这种精化意义的拓扑镜像对称，对作为加权射影空间超曲面得到的 95 个 K3 曲面而言已在 [Bel02] 中详细研究：其中 57 个在 95 个之中具有拓扑镜像，其余 38 个则不具有，这一点已为人知。
% [Tutor Note: 原文「残りの 38 個は持たない無い事が知られている」中「無い」为衍字，按原意译出。]

$n=3$ 时的权已由 [KS92a] 分类，已知共有 7555 个。朴素地考虑，这样得到的 Calabi--Yau 流形似乎会和 $\PP^4$ 的 5 次超曲面的情形一样，具有大的 $h^{1,n-1}$ 与小的 $h^{1,1}$，但照原样仍带有商奇点，因此必须对它作奇点解消。对个别例子具体构造解消未必显然，然而 Vafa [Vaf89] 在不具体构造奇点解消的情况下，导出了给出 Hodge 数生成函数 $P(t,\bar t):=\sum_{p,q=0}^{n}h^{p,n-q}t^p\bar t^q$ 的公式
\begin{equation}
  P(t,\bar t)=\Bigg(\sum_{l=0}^{h-1}\prod_{la_i/h\in\Z}\frac{1-(t\bar t)^{1-a_i/h}}{1-(t\bar t)^{a_i/h}}\prod_{la_i/h\notin\Z}(t\bar t)^{1/2-a_i/h}\left(\frac{t}{\bar t}\right)^{la_i/h-[la_i/h]-1/2}\Bigg)\Bigg|_{\mathrm{int}}
  \tag{9.26}
\end{equation}
其中 $[x]$ 表示不超过 $x$ 的最大整数，$(f)|_{\mathrm{int}}$ 表示从 $f$ 中去掉 $t$ 与 $\bar t$ 的分数幂所得的式子。
% [Tutor Note: (9.26) 中两个 Π 的指标 i 未写出，按惯例理解为对 i=1,…,n+2 取乘积；此处按原文保留。]

例如，对 $\PP^4$ 的 5 次超曲面的情形，注意到
\begin{align}
  \left(\frac{1-t^4}{1-t}\right)^5 &= (1+t+t^2+t^3)^5 \tag{9.27}\\
  &= 1+5t+15t^2+35t^3+65t^4+101t^5+135t^6 \tag{9.28}\\
  &\quad +155t^7+155t^8+135t^9+101t^{10}+65t^{11}+35t^{12} \tag{9.29}\\
  &\quad +15t^{13}+5t^{14}+t^{15} \tag{9.30}
\end{align}
(9.26) 中 $l=0$ 的项为
\begin{align}
  \left.\left(\frac{1-(t\bar t)^{1-1/5}}{1-(t\bar t)^{1/5}}\right)^5\right|_{\mathrm{int}} &= \left.\left(1+(t\bar t)^{1/5}+(t\bar t)^{2/5}+(t\bar t)^{3/5}\right)^5\right|_{\mathrm{int}} \tag{9.31}\\
  &= 1+101(t\bar t)+101(t\bar t)^2+(t\bar t)^3 \tag{9.32}
\end{align}
而 $l\ne 0$ 的项为
\begin{align}
  \sum_{l=1}^{4}\left[(t\bar t)^{\frac{1}{2}-\frac{1}{5}}\left(\frac{t}{\bar t}\right)^{\frac{1}{5}l-\left[\frac{1}{5}l\right]-\frac{1}{2}}\right]^5 &= \sum_{l=1}^{4}(t\bar t)^{\frac{5}{2}-1}\left(\frac{t}{\bar t}\right)^{l-\frac{5}{2}} \tag{9.33}\\
  &= \sum_{l=1}^{4}t^{l-1}\bar t^{4-l} \tag{9.34}\\
  &= \sum_{p=0}^{3}t^p\bar t^{3-p} \tag{9.35}
\end{align}

同样，对 $\PP^3$ 的 4 次超曲面（这是 K3 曲面）用 (9.26) 计算 Hodge 数，则由
\begin{equation}
  \left(\frac{1-t^3}{1-t}\right)^4=1+4t+10t^2+16t^3+19t^4+16t^5+10t^6+4t^7+t^8
  \tag{9.36}
\end{equation}
可知 (9.26) 中 $l=0$ 的项为
\begin{align}
  \left.\left(\frac{1-(t\bar t)^{1-1/4}}{1-(t\bar t)^{1/4}}\right)^4\right|_{\mathrm{int}} &= \left.\left(1+(t\bar t)^{1/4}+(t\bar t)^{2/4}\right)^4\right|_{\mathrm{int}} \tag{9.37}\\
  &= 1+19(t\bar t)+(t\bar t)^2 \tag{9.38}
\end{align}
% CHUNK c05 DONE pages p016–p021
